\chapter{\texorpdfstring{$\alpha$}{α} 消費関数}\label{ch:m09}

\begin{priobox}{C}
出題は\kakomon{2017}{2}の1回．$B(t)=\sqrt tZ(t)$ の平均・共分散と $\alpha(t)$ は与えられるが，
「中間で $\alpha(t_1)$ を消費し，最終では $\Prob(Z(t_1)\le z_1,\,Z(1)>z_2)=\alpha(1)-\alpha(t_1)$ で $z_2$ を決める」という原理は知らないと〔4〕が書けない．
本編の第14章14.2節．
\end{priobox}

% ============================================================
\section{概要}\label{sec:m09-gaiyo}
中間解析で有効性が示されれば試験を早く止めたいが，同じ $\alpha$ で何度も検定すると第1種の過誤確率が増える．
全体の $\alpha$ を各解析に「配分」するのが群逐次デザインで，その配分を情報時間 $t$ の関数で決めたものが $\alpha$ 消費関数である．
\begin{pointbox}[覚える式]
\begin{itemize}
\item 情報時間 $t=I/I_{\max}$．ログランク検定なら $t=$（その時点のイベント数）/（最終の目標イベント数）．
\item $B(t)=\sqrt t\,Z(t)$ は $H_0$ の下で $\E[B(t)]=0$，$\Cov[B(s),B(t)]=\min(s,t)$（近似的に標準 Brown 運動）．
\item $s<t$ なら $\Corr[Z(s),Z(t)]=\Cov[Z(s),Z(t)]=\sqrt{s/t}$．
\item $\alpha$ 消費関数 $\alpha(t)$：$\alpha(0)=0$，$\alpha(1)=\alpha$ の非減少関数．
      第 $k$ 回の解析で消費するのは $\alpha(t_k)-\alpha(t_{k-1})$．
\item 第1回：$\Prob(Z(t_1)>c_1)=\alpha(t_1)$ より $c_1=z_{\alpha(t_1)}$．
\item 第2回：$\Prob\bigl(Z(t_1)\le c_1,\ Z(t_2)>c_2\bigr)=\displaystyle\int_{-\infty}^{c_1}\int_{c_2}^{\infty}\phi_2\bigl(x_1,x_2;\sqrt{t_1/t_2}\bigr)\dd x_2\dd x_1=\alpha(t_2)-\alpha(t_1)$．
\end{itemize}
\end{pointbox}

% ============================================================
\section{詳細}\label{sec:m09-shosai}
\subsection{検定統計量の相関}
情報時間 $t$ までのデータによるスコア統計量を $S(t)$ とする（ログランクなら $U=\sum(d_{1j}-E_{1j})$）．$H_0$ の下で
(i) $\E[S(t)]=0$，$\V[S(t)]\approx tI_{\max}$，(ii) $s<t$ なら増分 $S(t)-S(s)$ は $S(s)$ と無相関（第2章の「分散を足せる理由」と同じ）．すると
\[
\Cov[S(s),S(t)]=\V[S(s)]+\Cov[S(s),S(t)-S(s)]=sI_{\max}.
\]
$Z(t)=S(t)/\sqrt{tI_{\max}}$，$B(t)=\sqrt tZ(t)=S(t)/\sqrt{I_{\max}}$ とおくと $\Cov[B(s),B(t)]=s=\min(s,t)$ で，
\[
\Cov[Z(s),Z(t)]=\frac{\Cov[B(s),B(t)]}{\sqrt s\sqrt t}=\frac{s}{\sqrt{st}}=\sqrt{\frac st}.
\]
$\V[Z(t)]=1$ なので，これは相関係数でもある．「$Z(t)$ は $Z(s)$ の情報に新しい情報を足したもので，重なりの割合が相関」と覚える．

\subsection{繰り返し検定の問題}
$t=0.5$ と $t=1$ の2回，それぞれ片側 $z_{0.025}=1.96$ で検定すると，全体の第1種の過誤確率は $1-\Prob(Z(0.5)\le1.96,\,Z(1)\le1.96)$（相関 $\sqrt{0.5}$）で，数値計算すると約 $0.042$ と名目の $0.025$ を大きく超える．

\subsection{限界値の決め方}
各回の事象「それまで棄却されず，第 $k$ 回で初めて棄却」は互いに排反なので，それぞれの確率を $\alpha(t_k)-\alpha(t_{k-1})$ にすれば，全体の第1種の過誤確率は足して $\alpha(t_K)=\alpha$．
\begin{enumerate}
\item $c_1$ は正規分布表で求まる．
\item $c_2$ 以降は多変量正規分布の数値積分が必要で，試験では「$\phi_2$ を用いた式で説明せよ」の形で問われる．
\end{enumerate}
$\alpha$ 消費関数法の利点は，解析の回数・時期を前もって固定しなくてよいことである．途中で解析を追加しても，消費する $\alpha$ は $\alpha(t)$ から自動的に決まる．

\subsection{代表的な消費関数}
\begin{center}\small
\begin{tabular}{@{}llL{70mm}@{}}
\toprule
名称 & $\alpha(t)$ & 特徴\\
\midrule
O'Brien--Fleming 型 & $2-2\Phi(z_{\alpha/2}/\sqrt t)$ & 初期はほとんど消費しない．最終の限界値が名目値に近い\\
Pocock 型 & $\alpha\log\{1+(e-1)t\}$ & 各回ほぼ同じ限界値．早期中止しやすいが最終が厳しい\\
べき型 & $\alpha t^\rho$ & $\rho$ が大きいほど初期の消費が少ない（2017年は $\rho=2$）\\
\bottomrule
\end{tabular}
\end{center}

\begin{warnbox}
\begin{itemize}
\item 最終解析の限界値を「$\Prob(Z(1)>c_2)=\alpha-\alpha(t_1)$」として正規分布表で引く（中間で棄却されなかったという条件が抜けている）．
\item 相関を $t_1/t_2$ とする（正しくは $\sqrt{t_1/t_2}$）．$\Cov[B]$ と $\Cov[Z]$ を混同する．
\item 解析を追加したときに，追加した回の消費を $\alpha(t_{\text{新}})$ とする（正しくは前回までとの差 $\alpha(t_{\text{新}})-\alpha(t_{\text{前}})$）．
\end{itemize}
\end{warnbox}

% ============================================================
\section{出題例}\label{sec:m09-shutsudai}
\begin{shutsudai}{\kakomon{2017}{2}：ログランク検定の群逐次デザイン}
片側5\%，目標1000イベント，$L=400$ で中間解析（$t=0.4$），$\alpha(t)=0.05t^2$．
〔1〕中間で $0.05\times0.16=0.008$，最終で $0.042$．
〔2〕$\Prob(Z(0.4)\ge z_1)=0.008$ より $z_1=2.41$．
〔3〕$\Cov[Z(0.4),Z(1)]=\Cov[B(0.4),B(1)]/\sqrt{0.4}=0.4/\sqrt{0.4}=\sqrt{0.4}$．
〔4〕$\int_{-\infty}^{z_1}\int_{z_2}^{\infty}\phi_2(x_1,x_2;\sqrt{0.4})\dd x_2\dd x_1=0.042$ を満たす $z_2$．
〔5〕$L=700$ を追加：2回目 $0.05\times0.49-0.008=0.0165$，最終 $0.05-0.0245=0.0255$．
\end{shutsudai}

% ============================================================
\section{例題}\label{sec:m09-reidai}
\begin{reidai}[\kakomon{2017}{2}型：$\alpha$ 消費関数による中間解析]\label{reidai:m09-1}
生存期間を評価項目とする2群のランダム化比較試験で，片側有意水準2.5\%のログランク検定（検定統計量が限界値より大きいとき棄却）を行う．
最終解析の目標イベント数は両群合わせて600，中間解析を1回，イベント数 $L$ の時点で行い，情報時間を $t=L/600$ とする．
情報時間 $t$ の検定統計量 $Z(t)$ は $H_0$ の下で標準正規分布に従い，$B(t)=\sqrt tZ(t)$ は $\E[B(t)]=0$，$\Cov[B(t_1),B(t_2)]=\min(t_1,t_2)$ を満たす．
$\alpha$ 消費関数を $\alpha(t)=0.025t^2$ とする．
\begin{enumerate}[label=〔\arabic*〕]
\item $L=300$ で中間解析を行うとき，中間解析と最終解析で消費する $\alpha$ を求めよ．
\item 中間解析の棄却限界値 $z_1$ を求めよ．
\item $\Cov[Z(0.5),Z(1)]$ を求めよ．
\item 最終解析の棄却限界値 $z_2$ の求め方を，平均 $\bm 0$・分散1・相関 $\rho$ の2変量正規密度 $\phi_2(x_1,x_2;\rho)$ を用いて説明せよ．
\item 計画を変更し，$L=480$ でも中間解析を行うことにした（合計3回）．2回目の中間解析と最終解析で消費する $\alpha$ を求めよ．
\end{enumerate}
\end{reidai}

% ============================================================
\section{例題の詳しい解説}\label{sec:m09-kaisetsu}
\begin{kaitou}{reidai:m09-1}
〔1〕$t_1=300/600=0.5$．中間解析で消費する $\alpha$ は $\alpha(0.5)=0.025\times0.25=0.00625$．
最終解析で消費するのは $\alpha(1)-\alpha(0.5)=0.025-0.00625=0.01875$．

〔2〕$\Prob(Z(0.5)>z_1\mid H_0)=0.00625$ を満たす $z_1$．正規分布表で上側確率0.0064（2.49）と0.0062（2.50）の間なので $z_1\approx2.50$．
（名目の $z_{0.025}=1.96$ よりかなり厳しい．）

〔3〕$Z(0.5)=B(0.5)/\sqrt{0.5}$，$Z(1)=B(1)$ なので
\[
\Cov[Z(0.5),Z(1)]=\frac{1}{\sqrt{0.5}}\Cov[B(0.5),B(1)]=\frac{\min(0.5,1)}{\sqrt{0.5}}=\frac{0.5}{\sqrt{0.5}}=\sqrt{0.5}=0.707 .
\]
$\V[Z(t)]=1$ なのでこれは相関係数でもある．

〔4〕$(Z(0.5),Z(1))$ は $H_0$ の下で平均 $\bm0$，分散1，相関 $\sqrt{0.5}$ の2変量正規分布に従う（近似）．
最終解析で棄却するのは「中間で棄却されず（$Z(0.5)\le z_1$），最終で $Z(1)>z_2$」のときなので，$z_2$ は
\[
\Prob\bigl(Z(0.5)\le z_1,\ Z(1)>z_2\bigr)=\int_{-\infty}^{z_1}\int_{z_2}^{\infty}\phi_2\bigl(x_1,x_2;\sqrt{0.5}\bigr)\dd x_2\,\dd x_1=0.01875
\]
を満たす値として，数値積分で求める．
全体の第1種の過誤確率は，排反な2つの事象の確率の和 $0.00625+0.01875=0.025$ になる．
（参考：数値計算すると $z_2\approx2.02$．中間解析をしない場合の $1.96$ よりわずかに大きい．）

〔5〕$t_2=480/600=0.8$．$\alpha(0.8)=0.025\times0.64=0.016$．
2回目の中間解析で消費するのは $\alpha(0.8)-\alpha(0.5)=0.016-0.00625=0.00975$，
最終解析で消費するのは $\alpha(1)-\alpha(0.8)=0.025-0.016=0.009$．
（合計 $0.00625+0.00975+0.009=0.025$．）
\end{kaitou}

\begin{summarybox}
\begin{itemize}
\item 情報時間 $t$＝イベント数の割合．$\Cov[B(s),B(t)]=\min(s,t)$，$\Corr[Z(s),Z(t)]=\sqrt{s/t}$．
\item 第 $k$ 回の消費 $=\alpha(t_k)-\alpha(t_{k-1})$．$c_1=z_{\alpha(t_1)}$，$c_2$ は「それまで棄却されず第2回で棄却」の確率で決める（2変量正規の積分）．
\item 解析を追加しても消費は $\alpha(t)$ から自動的に決まる．
\end{itemize}
\end{summarybox}
